\documentclass{article}
\usepackage{amsmath, amsthm, amsfonts, amssymb}
\theoremstyle{plain}
\newtheorem{thm}{Theorem}
\numberwithin{thm}{section}

\theoremstyle{plain}
\newtheorem{prop}{Proposition}
\numberwithin{prop}{section}

\theoremstyle{definition}
\newtheorem{defn}{Definition}
\numberwithin{defn}{section}

\theoremstyle{remark}
\newtheorem*{rem}{Remark}

\newtheorem*{cor}{Corollary}

\numberwithin{equation}{section}

\newcommand{\op}{\prime}
\DeclareMathOperator{\re}{Re}
\DeclareMathOperator{\im}{Im}
\DeclareMathOperator{\curl}{curl}
\begin{document}
\title{Sequences and Series}
\author{Amey Joshi}
\maketitle
\section{Important inequalities}\label{s1}
\begin{prop}[Triangle inequality]\label{s1p1}
If $z_1, z_2 \in \mathbb{C}$ then $|z_1 + z_2| \le |z_1| + |z_2|$.
\end{prop}
\begin{proof}
$|z_1 + z_2|^2 = |z_1|^2 + |z_2|^2 + \bar{z}_1z_2 + z_1\bar{z_2} = |z_1|^2 +
|z_2|^2 + 2\re(\bar{z}_1z_2) \le |z_1|^2 + |z_2|^2 + 2|z_1||z_2|$, from which
the inequality follows immediately.
\end{proof}

\begin{prop}[Reverse triangle inequality]\label{s1p2}
If $z_1, z_2 \in \mathbb{C}$ then $||z_1| - |z_2|| \le |z_1 - z_2|$.
\end{prop}
\begin{proof}
$|z_1| = |z_1 - z_2 + z_2| \le |z_1 - z_2| \le |z_2|$ so that $|z_1| - |z_2| 
\le |z_1 - z_2|$. Similarly, $|z_2| - |z_1| \le |z_1 - z_2|$, so that $||z_1| - 
|z_2|| \le |z_1 - z_2|$.
\end{proof}

\begin{prop}[Cauchy-Schwarz inequality]\label{s1p3}
If $\{a_n\}$ and $\{b_n\}$ are sequences of complex numbers then
\[
\Big|\sum_{n=1}^Na_n\bar{b}_n\Big|^2 \le \Big|\sum_{n=1}^Na_n\Big|^2 
\Big|\sum_{n=1}^Nb_n\Big|^2
\]
\end{prop}
\begin{proof}
Consider two $N$-dimensional complex vectors $a = (a_1, \ldots, a_N)$ and 
$b = (b_1, \ldots, b_N)$. Let $c = a - \lambda b$ for a complex number 
$\lambda$. The norm of a complex vector is defined as $\lVert a \rVert^2 = 
\sum_{k=1}^n |a_k|^2$. Clearly, $\lVert a \rVert^2 \ge 0$. Consider $\lVert c
\rVert^2 = \lVert a - \lambda b\rVert^2 = (a - \lambda b, \bar{a} - 
\bar{\lambda}\bar{b}) = \lVert a \rVert^2 - \bar{\lambda}(a, b) - \lambda(b, a)
+ \lVert b \rVert^2$, where the inner product $(a, b)$ is
\[
\sum_{k=1}^N a_k \bar{b}_k.
\]
Clearly, $(b, a) = \overline{(a, b)}$. Choose $\lambda = (a, b)/r$ so that
\[
\lVert a \rVert^2 + \lVert b \rVert^2 \ge \frac{2|(a, b)|^2}{r}
\]
Now choose 
\[
r = \frac{2\lVert a \rVert^2 \lVert b \rVert^2}{\lVert a \rVert^2 + \lVert b \rVert^2}
\]
so that
\[
\lVert a \rVert^2 \lVert b \rVert^2 \ge |(a, b)|^2.
\]
Expressing the norms and inner products in terms of components, we get
\[
\Big|\sum_{k=1}^N |a_k|^2\Big|\Big|\sum_{k=1}^N |b_k|^2\Big| \ge 
\Big|\sum_{k=1}^N a_k\bar{b}_k\Big|^2
\]
\begin{eqnarray*}
\end{eqnarray*}
\end{proof}

\section{Limits, continuity and derivatives}\label{s2}
\begin{defn}\label{s2d1}
Let $f: \mathbb{C} \mapsto \mathbb{C}$ be defined at least in the deleted
$\delta$-neighbourhood of $z_0 \in \mathbb{C}$. Then 
\[
\lim_{z \rightarrow z_0} = w_0
\]
if for a given $\epsilon > 0, \exists\;\delta > 0$ such that for all $|z - z_0| 
< \delta$, $|f(z) - w_0| < \epsilon$.
\end{defn}
\begin{rem}
The definition does not require $f$ to be defined at $z_0$. Furthermore, $|z -
z_0| < \delta$ allows $z$ to approach $z_0$ from any point on the circle 
defined by $|z - z_0| = \delta$.
\end{rem}

\begin{defn}\label{s2d2}
A function $f:A \mapsto \mathbb{C}$ is continuous at a point $z_0 \in A 
\subset \mathbb{C}$ if
\[
\lim_{z \rightarrow z_0} = f(z_0).
\]
Equivalently, for a given $\epsilon > 0, \exists\;\delta > 0$ such that 
$|z - z_0| < \delta \Rightarrow |f(z) - f(z_0)| < \epsilon$.
\end{defn}

A useful application of this definition is the fact
\begin{prop}\label{s2p1}
The modulus function is continuous.
\end{prop}
\begin{proof}
Fix an $\epsilon > 0$. Then $||z| - |z_0|| \le \epsilon$ if $|z - z_0| <
\epsilon$, where we used the reverse triangle inequality \ref{s1p2}.
\end{proof}

\begin{defn}\label{s2d3}
The derivative of a function $f: A \mapsto \mathbb{C}$ at $z_0 \in A \subset
\mathbb{C}$ is the limit,
\[
Df = \lim_{|h| \rightarrow 0}\frac{f(z_0 + h) - f(z_0)}{h}.
\]
The function $f$ is said to be differentiable at $z_0$. $f$ is differentiable
in $A$ if it is differentiable at all points in $A$.
\end{defn}
\begin{rem}
The point $z_0 + h$ could be any point on the circle centred at $z_0$ and with
radius $|h|$.
\end{rem}

\begin{prop}\label{s2p2}
If $f$ is differentiable at $z$ then it is continous at $z$.
\end{prop}
\begin{proof}
Denote by $L$ the derivative of $f$ at $z$. Then for a given $1 > \epsilon_1 > 0$,
there exists $\delta_1 > 0$ such that 
\[
|h| < \delta_1 \Rightarrow \Bigg|L - \frac{f(z+h) - f(z)}{h}\Bigg| < \epsilon_1.
\]
That is, $|Lh - f(z+h) + f(z)| < \epsilon_1 |h|$. Using the reverse triangle
inequality, we get
\[
||Lh| - |f(z+h) - f(z)|| < \epsilon_1 |h| \Rightarrow -\epsilon_1 |h| < 
|f(z+h)-f(z)| - |Lh| < \epsilon_1 |h| 
\]
or,
\[
(|L| - \epsilon_1)|h| < |f(z+h) - f(z)| < (|L| + \epsilon_1)|h| < (|L| + 1)|h|.
\]
Thus, we have $|h| < \delta_1 \Rightarrow |f(z+h) - f(z)| < (|L| + 1)|h| <
(|L| + 1)\delta_1$. For a given $\epsilon > 0$, choose 
\[
\delta = \min\left(\delta_1, \frac{\epsilon}{(|L|+1)}\right)
\]
so that $|h| < \delta \Rightarrow |f(z+h) - f(h)| < \epsilon$.
given $\epsilon$, choose
\end{proof}

\begin{prop}\label{s2p3}
If $f(z) = u(x, y) + iv(x, y)$ is differentiable at a point $z$ then $u_x = v_y$
and $u_y = -v_x$, where the subscripts denote the partial derivatives.
\end{prop}
\begin{proof}
If we take the limit in the definition of the derivative along the real axis 
then $Df = u_x + iv_x$. If we take is along the imaginary axis then $Df = -iu_y
+ v_y$. Since the limit taken from all directions is identical if it exists,
we have $u_x = v_y$ and $u_y = -v_x$.
\end{proof}

\begin{rem}
The relations $u_x = v_y$ and $u_y = -v_x$ on the real and imaginary parts of
a differentiable function are called the \emph{Cauchy-Riemann} conditions.
\end{rem}

An immediate conseqence of the Cauchy-Riemann conditions is 
\begin{prop}\label{s2p4}
$u$ and $v$ are harmonic, that is, they satisfy Laplace's equation $\Delta u=0$
and $\Delta v = 0$.
\end{prop}

We will next prove the converse of proposition \ref{s2p3}, that is,
\begin{prop}\label{s2p5}
If $f(z) = u(x, y) + iv(x, y)$, $u_x = v_y$ and $u_y = -v_x$ at a point
$z = x + iy$ and the four partial derivatives are continuous then $f$ is 
differentiable at $z$.
\end{prop}
\begin{proof}
Let $h = k + il$. Since $u$ and $v$ are differentiable, for $|h| \rightarrow 0$,
we have
\begin{eqnarray*}
u(x + k, y + l) &=& u(x, y) + ku_x + lu_y \\
v(x + k, y + l) &=& v(x, y) + kv_x + lv_y
\end{eqnarray*}
so that $f(z + h) - f(z) = k(u_x + iv_x) + l(u_y + iv_)$ and
\[
\frac{f(z + h)-f(z)}{h} = \frac{(k - il)k(u_x + iv_x) + l(u_y + iv_)}{k^2+l^2}.
\]
Using $u_x = v_y$ and $u_y = -v_x$, we get
\[
\frac{f(z + h) - f(z)}{h} = \frac{(k^2 + l^2)(u_x + iv_x)}{k^2+l^2} =
u_x + iv_y.
\]
as $|h| \rightarrow 0$.
\end{proof}

From propositions \ref{s2p3} and \ref{s2p5}, we have
\begin{thm}\label{s2t1}
The necessary and sufficient conditions for $f(z) = u(x, y) + iv(x, y)$ to be
differentiable at a point is that the Cauchy-Riemann conditions are true at
that point and the four partial derivatives are continuous at it.
\end{thm}

\begin{defn}
A function $f$ is said to be holomorphic or analytic at a point $z_0$ if it is 
differentiable in a neighbourhood of $z_0$. $f$ is holomorphic in a domain $D$
if it is holomorphic at all points in $D$.
\end{defn}
Recall that a domain is an open, connected set. Therefore, $f$ is holomorphic
in $D$ means that for every $z_0 \in D$, we can find a neighbourhood $|z - z_0|
< \epsilon$ in which $f$ is differentiable. The term \emph{analytic} is also
used in the context of real-valued functions. However, it means that the 
function can be expressed as a power series in the neighbourhood of that point.
A real-valued function can be in $C^\infty$ and yet not be analytic. For example,
$\exp(-1/x^2)$ has all derivatives but it does not have a power series 
representation in the neighbourhood of $x = 0$ because all derivatives vanish.
Because differentiability and existence of a power series are not synonymous
outside complex-valued functions, we will prefer the term holomorphic which 
implies synonymity over analytic which does not.

\begin{defn}\label{s2d5}
A point at which a function fails to be analytic is called its singular point.
\end{defn}

\begin{defn}\label{s2d6}
A function analytic at all points in the finite complex plane is called entire.
\end{defn}

\section{Multi-valued functions}\label{s3}
Consider the function $f(z) = \log z$. If $z = re^{i\theta}$ then adding an
integral multiple of $2\pi$ to $\theta$ brings us to the same point. However,
$\log(re^{i(\theta + 2\pi m)}) = \log r + i\theta + 2\pi i m$. Thus, each time
one goes around the origin along a circle $|z| = r$, one comes back to same $z$
but has a different value of $\log z$. The point $z = 0$ is called a 
\emph{branch point} of the function $\log$. The function $\log$ fails to be
analytic unless we stay on a single branch of the function specified by a 
particular value of $m$. The branch $m = 0$ is called its \emph{principal
branch} and the function $\log$ is said to be \emph{multi-valued}.

The function $f(z) = \sqrt{z}$ is also multi-valued and has $z = 0$ as its
branch point because if $z = r\exp(i(\theta + 2\pi m))$ then $\sqrt{z} = 
\sqrt{r}\exp(i(\theta/2 + \pi m))$. If $m$ is even, we stay on one branch but
move to another one if $m$ is odd. Once again, the principal branch is the one
on which $m = 0$.

It is easy to see that the branch point of $\log(z - z_0)$ and $\sqrt{z - z_0}$
is $z = z_0$ and not the origin. The points $z_0 + re^{i(\theta + 2\pi m)}$
reveals the multi-valuedness.

A function remains single-valued if we go aroung a branch point as long as we
do not cross a \emph{branch cut}. A branch cut is an arbitrarily chosen line 
used to mark the boundary of different branches of a multi-valued function.
It is chosen out of convenience. For example:
\begin{enumerate}
\item The branch cut of $\log$ is often taken to be the positive real axis.
\item The branch cut of $\sqrt{z}$ is also taken to be the positive real axis.
\item The function $f(z) = \sqrt{(z - a)(z - b)}$, where $a < b \in \mathbb{R}$
is taken as the segment $[a, b]$ or the disjoint segments $[\infty, a] \cup [b,
\infty]$. Note that in the complex plane we have only one infinity, which is
mapped to the north pole of the Riemann sphere.
\end{enumerate}

Multi-valued functions are discontinuous across a branch cut. As a result, they
are not differentiable and hence not holomorphic.

\section{Integration}\label{s4}
Integration of a complex-valued function along a contour $C$ in the complex
plane is defined as the usual Riemann line integral. That is,
\[
\int_C f(z)dz = \lim_{n \rightarrow \infty}\sum_{k=1}^n f(z_k)\delta z_k,
\]
where the contour $C$ is divded into linear segments at points $z_1, \ldots,
z_n$ and $\delta z_k = |z_{k+1} - z_k|$. The integral is said to exist when the 
limit does. Using this definition, we can easily get an estimate of the 
magnitude of the integral in the following proposition.

\begin{prop}[The ML inequality]\label{s4p1}
If $|f(z)| \le M$ for all $z$ on a contour of length $L$ then 
\[
\int_C f(z)dz \le ML.
\]
\end{prop}
\begin{proof}
Using the definition of the  Riemann integral
\[
\int_C f(z)dz = \lim_{N \rightarrow \infty}\sum_{k=1}^Nf(z_k)\delta z_k,
\]
where we split the contour $C$ into segments joining the points $\{z_k\}$ with
$k = 1, \ldots, N$. We denote the length of the segment joining $z_k$ and 
$z_{k+1}$ by $\delta z_k$. Thus,
\[
\Big|\int_C f(z)dz\Big| = 
\lim_{N \rightarrow \infty}\Big|\sum_{k=1}^N f(z_k)\delta z_k\Big|,
\]
where we used the continuity of the modulus function to interchange the limit
and the modulus operations. Using the triangle inequality of \ref{s1p1}, we get
\[
\Big|\int_C f(z)dz\Big| \le 
\lim_{N \rightarrow \infty}\sum_{k=1}^N|f(z_k)|\delta z_k
\le \lim_{N \rightarrow \infty}M \sum_{k=1}^N |\delta z_k| = ML,
\]
where we used the fact that
\[
\lim_{N \rightarrow \infty}\sum_{k=1}^N |\delta z_k| = L.
\]
\end{proof}

Green's theorem in vector calculus says that
\[
\oint_C \vec{u}.\cdot d\vec{x} = \iint_S \curl\vec{u}\cdot d\vec{a},
\]
where $C$ is a closed curve enclosing an area $S$, $d\vec{x}$ is the line 
element and $d\vec{a}$ the area element. If $\vec{u} = u\vec{e}_x + v\vec{e}_y$
then 
\[
\curl\vec{u} = \left(\frac{\partial v}{\partial x} - 
\frac{\partial u}{\partial y}\right)\vec{e}_z
\]
and Green's theorem becomes
\[
\oint_C (udx + vdy) = \iint_S\left(\frac{\partial v}{\partial x} - 
\frac{\partial u}{\partial y}\right)dxdy.
\]
We will apply this to the contour integral of a complex-valued function and
prove an important theorem.
\begin{thm}[Cauchy]\label{c4t1}
If $f$ is holomorphic on and inside a simple closed curve $C$ then
\[
\oint_C f(z)dz = 0.
\]
\end{thm}
\begin{proof}
Let $f(z) = u(x, y) + iv(x, y)$ and $dz = dx + idy$. Then $fdz = (udx - vdy) +
i(udy + vdx)$. Applying Green's theorem in two dimensions to the integrals on 
the right,
\[
\oint_C (udx - vdy) = \iint_S\left(-\frac{\partial v}{\partial x} - 
\frac{\partial u}{\partial y}\right)dxdy = 0,
\]
using the Cauchy-Riemann condtion. Similarly,
\[
\oint_C (vdx + udy) = \iint_S\left(\frac{\partial u}{\partial x} - 
\frac{\partial v}{\partial y}\right)dxdy = 0,
\]
once again due to Cauchy-Riemann condition. Thus,
\[
\oint_C f(z)dz = 0.
\]
\end{proof}

\begin{thm}[Cauchy integral formula]\label{s4t2}
Let $f$ be holomorphic on and inside a simple, closed contour $C$. Then for any
interior point $z_0$,
\[
f(z_0) = \frac{1}{2\pi i}\oint_C \frac{f(z)}{z - z_0}dz.
\]
\end{thm}
\begin{proof}
Modify the contour $C$ to create a key-hole $C_r$ of small radius $r$ around the 
point $z_0$. In the interior of the modified contour, the integrand is 
holomorphic and therefore the contour integral is zero. The integral over the
modified contour can be written as a sum of the integral along $C$ traversed
anti clockwise and that round $C_r$ traversed clock-wise. Thus, if $C^\op$ is
the modified contour,
\[
\oint_{C^\op} = \oint_C - \oint_{C_r}
\]
As the lhs is zero, we have
\[
\oint_C \frac{f(z)}{z - z_0}dz = \oint_{C_r}\frac{f(z)}{z - z_0}dz.
\]
Along $C_r$, $z = z_0 + r\exp(i\theta)$ so that $dz = ir\exp(i\theta)$ and
the integral is from $0$ to $2\pi$. Thus,
\[
\oint_C \frac{f(z)}{z - z_0}dz = i\int_0^{2\pi}f(z_0 + re^{i\theta})d\theta.
\]
If we now let $r \rightarrow 0$, the rhs evaluates to $2\pi if(z_0)$.
\end{proof}

\begin{thm}\label{s4t3}
If $f$ is holomorphic on and inside a simple closed contour $C$ it has 
derivatives of all orders inside $C$. Furthermore,
\[
(D^k f)(z) = \frac{k!}{2\pi i}\oint_C\frac{f(\zeta)}{(\zeta - z)^{k+1}}d\zeta.
\]
\end{thm}
\begin{proof}
From theorem \ref{s4t2},
\[
f(z) = \frac{1}{2\pi i}\oint_C \frac{f(\zeta)}{\zeta - z}d\zeta.
\]
so that
\[
f(z + h) = \frac{1}{2\pi i}\oint_C \frac{f(\zeta)}{\zeta - z - h}d\zeta.
\]
and
\[
\frac{f(z + h) - f(z)}{h} = \frac{1}{2\pi i}\oint_C \frac{f(\zeta)}{(\zeta - z)
(\zeta - z - h)} d\zeta.
\]
Upon taking the limit $h \rightarrow 0$, we get
\[
(Df)(z) = \frac{1}{2\pi i}\oint_C \frac{f(\zeta)}{(\zeta - z)^2}d\zeta.
\]
If we assume that
\[
(D^k f)(z) = \frac{k!}{2\pi i}\oint_C \frac{f(\zeta)}{(\zeta - z)^{k+1}}d\zeta
\]
then we the expression for $(Df)(z)$ serves as the base case and by substituting
\[
g_k = \frac{f(\zeta)}{(\zeta - z)^k}
\]
we can show that
\[
(Dg_k)(z) = \frac{1}{2\pi i}\oint_C \frac{g_k(\zeta)}{(\zeta - z)^2}d\zeta
\]
which proves the induction hypothesis. 
\end{proof}

Being holomorphic is a very stringent condition on a function. It ensures that
existence of the first derivative guarantees the existence of all higher order
derivatives.

Theorem \ref{s4t3} gives us an expression for the $k$th derivative of a 
holomorphic function as an integral. This allows us to get an estiamte of its
magnitude.
\begin{prop}\label{s4p2}
If $f$ is holomorphic on and inside a simple, closed contour $C$ then
\[
|(D^n f)(z)| \le \frac{n! M}{R^n}
\]
where $M$ is an upper bound of $|f|$ and $R$ is an upper bound of $|\zeta - z|$.
\end{prop}
\begin{proof}
From the expression of the $n$-th derivative,
\[
|(D^n f)(z)| \le 
\frac{n!}{2\pi}M\Big|\oint_C\frac{d\zeta}{(\zeta - z)^{n+1}}\Big|
\]
If $R > |\zeta - z|$ then
\[
|(D^n f)(z)| \le \frac{Mn!}{2\pi R^{n+1}}\Big|\oint_C d\zeta\Big|
\]
Since $R$ is an upper bound to $|\zeta - z|$, $C$ is enclosed by a circle of
radius $R$ centred at $z$. Therefore, the length of the contour are at most
$2\pi R$. Therefore,
\[
|(D^n f)(z)| \le \frac{Mn!}{R^n}
\]
\end{proof}

This proposition has a severe consequence.
\begin{thm}[Liouville]\label{s4t4}
If $f$ is entire and bounded then it is a constant.
\end{thm}
\begin{proof}
If $f$ is entire then the magnitude of its first derivative is bounded by
\[
|Df(z)| \le \frac{M}{R}
\]
where $R$ can be allowed to take as large a value as we wish so that $Df = 0$
and hence $f$ is a constant.
\end{proof}

We can use Liouville's theorem to prove
\begin{thm}[Fundamental theorem of algebra]\label{s4t5}
A polynomial in $z$ of a finite degree $n > 0$ and constant coefficients has 
$n$ roots.
\end{thm}
\begin{proof}
Let $P$ be a polynomial in $z$ with constant coefficients. Let, if possible, $f$
have no zeros. Therefore the function $f = 1/P$ is entire. Furthermore, as $|z|
\rightarrow \infty$, $|P| \rightarrow \infty$ if $n > 0$. Therefore 
\[
\lim_{|z| \rightarrow \infty} f(z) = 0.
\]
Therefore, by Liouville's theorem $f$ must be a constant equal to zero although
$P$ has no zeros. This is a contradiction. Therefore, $P$ must have at least
one zero, say $z_1$. Now write $P = (z - z_1)P_1$, where $P_1$ is a polynomial 
of degree $n - 1$. Conside $f_1 = 1/P_1$. If $P_1$ has no zeros, we once again
arrive at the contradiction that $f_1 = 0$. Therefore, $P_1$ too has at least
one root, say $z_2$. If we continue this way, in the $k$th step, we are left
with a polynomial of degree $n - k$. Since $n$ is finite, this procedure 
terminates in $n$ steps giving $n$ roots.
\end{proof}

It is important to note that the degree of a polynomial is finite. This allows
us to conclude that $|P| \rightarrow \infty$ as $|z| \rightarrow \infty$. If
we decide to treat the power series of $e^z$ as a polynomial, we cannot say
that because we know that $|e^z| \rightarrow 0$ as $\re(z) \rightarrow -\infty$.

\section{Singularities}\label{s5}
A complex-valued function can fail to be analytic at a point for several reasons.
In this section, we will examine a few of them. Each one reslts in a different
kind of a singularity.


\section{Sequences and series}\label{s6}

\section{Interesting theorems and problems}\label{s1a}
\section{Theorems}
\begin{thm}\label{s1at1}
Let $f$ be a meromorphic function with $N$ simple zeros and $M$ simple poles
inside a circle $C$. Then
\[
\frac{1}{2\pi i}\oint_C \frac{Df}{f}dz = N - M.
\]
\end{thm}
\begin{proof}
If $a$ is a simple zero of $f$ then $f = (z - a)g(z)$ where $g(a) \ne 0$. Its
derivative is $Df = g(z) + (z - a)Dg$ so that
\[
\frac{Df}{f} = \frac{g}{f} + \frac{Dg}{g} = \frac{1}{z - a} + \frac{Dg}{g}.
\]
If $C_a$ is a small enough contour centred at $z = a$ which encloses no other
pole or zero of $f$,
\[
\oint_{C_a}\frac{Df}{f}dz = 2\pi i.
\]
If $p$ is a simple pole of $f$ then $f = h(z)/(z - p)$ and 
\[
Df = \frac{Dh}{z - p} - \frac{h}{(z - p)^2}
\]
so that
\[
\frac{Df}{f} = \frac{Dh}{h} - \frac{1}{z - p}
\]
and hence
\[
\oint_{C_p}\frac{Df}{f}dz = -2\pi i,
\]
where $C_p$ is a small enough contour centred at $z = p$ that encloses no other
pole or zero of $f$.

If we now distort $C$ with a key hole for each zero and pole so that it encloses
only one at a time, each zero contributes $+2\pi i$ and each pole contributes
$-2\pi i$.
\end{proof}

\begin{thm}[Riemann theorem of removable singularity]\label{s1at2}
Let $f$ be a bounded, holomorphic function on a punctured disc $0 < |z - z_0|
< \delta$ the $z_0$ is a removable singularity.
\end{thm}
\begin{proof}
Since $f$ is bounded and holomorphic on a puctured disc centered at $z_0$, it
has a Laurent series expansion of the form
\[
f(z) = \sum_{n = -\infty}^\infty a_n(z - z_0)^n.
\]
The coefficients $a_n$ are obtained using the formula
\[
a_n = \frac{1}{2\pi i}\oint_C \frac{f(z)}{(z - z_0)^{n + 1}}dz,
\]
where $C$ is a simple, closed curve inside the punctured disc. For a fixed
$k > 0$,
\[
a_{-k} = \frac{1}{2\pi i}\oint_C f(z)(z - z_0)^{k - 1}dz
\]
so that if $C$ is a circle of radius $0 < r < \delta$,
\[
|a_{-k}| \le \frac{1}{2\pi}Mr^{k - 1}2\pi r = Mr^k,
\]
where $M$ is an upper bound of $f$ on the punctured disc. Now, $a_{-k}$ is
a constant and therefore should be independent of the contour's radius. The
only way that can happen is $a_{-k} = 0$ for all $k > 0$ so that the Laurent
series becomes a Taylor series, making $z = z_0$ a removable singularity.
\end{proof}

\begin{thm}[Casorati-Weierstrass]\label{s1at3}
Let $z_0$ be an essential singularity of a function $f$. Define the punctured
disc centred at $z_0$ as $D^\op = \{z \in \mathbb{C}\;:\; 0 < |z - z_0| < 
\delta\}$. For any $\delta > 0$, howsoever small, the image of $D^\op$ is dense
in $\mathbb{Z}$.
\end{thm}
\begin{proof}
Let, if possible, the image of $D^\op$ not be dense in $\mathbb{C}$. Therefore,
there is at least one point $w_0 \in \mathbb{Z}$ and an $\epsilon > 0$ such that
no point of $D^\op$ is mapped into $|w - w_0| < \epsilon$ under $f$. Define a
function $g: D^\op \rightarrow \mathbb{C}$ as 
\[
g(z) = \frac{1}{f(z) - w_0}.
\]
Since $z \in D^\op$, the function is defined for all points in $D^\op$. Since
$|f(z) - w_0| > \epsilon$ for all $z \in D^\op$, 
\[
|g(z)| = \frac{1}{|f(z) - w_0|} < \frac{1}{\epsilon}.
\]
Thus $g$ is bounded and entire in $D^\op$. By Riemann's theorem of removable
singularities, the point $z_0$ is a removable singularity of $g$. Therefore,
$g$ has a power series of the form
\[
g(z) = \sum_{k \ge 0}a_k(z - z_0)^k
\]
Since
\[
f(z) = w_0 + \frac{1}{g(z)}
\]
we have three possibilities:
\begin{enumerate}
\item If $g(z_0) \ne 0$, $z_0$ is a removable singularity of $f$.
\item If $g(z_0) = 0$ but not all its derivatives at $z = z_0$ are not zero.
If $k$ is the derivative of lowest order which does not vanish at $z = z_0$
then $z_0$ is a pole of order $k$ of the function $f$.
\item $g = 0$ identically for all $z \in D^\op$. This is impossible since 
$g = 1/(f - w_0)$.
\end{enumerate}
Since $z_0$ was an essential singularity the first two possibilities are 
contradictory. Therefore, our assumption that image of $f$ is not dense is false.
\end{proof}
\end{document}
